\chapter{Podsjetnik za dan obrane}
\label{ch:podsjetnik}

Ovo je zadnje što treba pročitati --- ujutro, prije izlaska. Sve je izvedeno u
prethodnim poglavljima; ovdje stoji samo ono što se pamti.

\section{Cijeli seminar na jednoj stranici}
\label{sec:pod-pregled}

\begin{longtable}{@{}p{0.09\linewidth} p{0.24\linewidth} p{0.28\linewidth} p{0.31\linewidth}@{}}
\toprule
& \textbf{Klasa} & \textbf{Ključna ideja} & \textbf{Rezultat} \\
\midrule
\endhead
\textbf{Z1} & QP & KKT uvjeti, $2^m$ kombinacija aktivnosti; dualnost i cijene u
sjeni & $\xopt = (190/33,\ 140/33)$, $\popt = 2000/33$, $\muopt = (404/33,\ 0)$ \\[6pt]
\textbf{Z2} & QP (robusni) & Supremum po pravokutniku postignut u vrhu $\Rightarrow$
osam nejednakosti & $x_{\mathrm{rob}}^{\star} = (6{,}5,\ 179/34)$,
$f = 84{,}624$, $+39{,}63\%$ \\[6pt]
\textbf{Z3} & LP & $\ell_1$-norma pomoćnim varijablama postaje linearna;
$6N$ nejednakosti & $J^{\star} = 1{,}9659$ (A) i $2{,}5060$ (B) pri $N = 100$ \\[6pt]
\textbf{Z4} & SDP & Zajednička kvadratna Ljapunovljeva funkcija; kongruencija
$\Rightarrow$ BMI postaje LMI & $K = [-3{,}4246,\ -0{,}1047]$,
$\alpha = 0{,}2884\ \mathrm{s}^{-1}$ \\
\bottomrule
\end{longtable}

\begin{pojam}{Hijerarhija klasa}
\begin{equation}
\text{LP} \subset \text{QP} \subset \text{SOCP} \subset \text{SDP}
\end{equation}
Zadatak 3 je LP, Zadaci 1 i 2 su QP, Zadatak 4 je SDP. Sva su četiri
\textbf{konveksna}, sva se rješavaju metodama unutarnje točke, i za sva vrijedi da
je svaki lokalni optimum ujedno globalni.
\end{pojam}

\section{Tri ideje koje se ponavljaju}
\label{sec:pod-ideje}

\begin{korak}[Ako zaboraviš sve ostalo, zapamti ovo troje]
\begin{enumerate}[leftmargin=*]
  \item \textbf{Reformulacija, ne novi algoritam.} Sva četiri zadatka rješavaju se
        tako da se problem \emph{prepiše} u standardnu klasu, a onda preda
        gotovom rješavaču. Apsolutna vrijednost $\to$ pomoćne varijable
        (Z3); umnožak nepoznanica $\to$ supstitucija $Y = KP^{-1}$ (Z4);
        beskonačno mnogo ograničenja $\to$ vrhovi (Z2).
  \item \textbf{Beskonačno $\to$ konačno preko konveksne ljuske.} Linearna
        funkcija po konveksnoj ljusci postiže ekstrem u generatoru. To vrijedi za
        vrhove pravokutnika u Z2 i za matrice modova u Z4 (Predavanje 7, slajd
        16/37). Ista rečenica, dvije primjene.
  \item \textbf{Množitelji imaju značenje.} $\mu_i^{\star}$ nije samo broj koji
        ispadne iz računa --- to je cijena $i$-tog ograničenja,
        $\partial\popt/\partial v_i = -\mu_i^{\star}$. Nulti množitelj znači
        neaktivno ograničenje, veliki množitelj znači ograničenje koje stvarno
        veže.
\end{enumerate}
\end{korak}

\section{Pet zamki koje se najskuplje plaćaju}
\label{sec:pod-zamke}

\begin{zamka}[1. ,,Oba moda stabilna, dakle sustav stabilan`` (Z4)]
Najteža greška u cijelom seminaru. Hurwitzovost svakog moda \emph{ne} povlači
stabilnost sklopnog sustava --- $\rho(\Phi_{\mathrm{cikl}}) = 9e^{-2} = 1{,}218 > 1$.
Pod 4a stani na ,,svaki mod zasebno``.
\end{zamka}

\begin{zamka}[2. ,,Funkcija cilja je konveksna`` kao cijeli odgovor (Z1)]
Pitanje je je li \emph{problem} konveksan. Odgovor mora pokriti i dopustivi skup,
i spomenuti egzistenciju (koercivnost, jer je skup neograničen) i jedinstvenost
(stroga konveksnost).
\end{zamka}

\begin{zamka}[3. Negativan množitelj shvaćen kao greška (Z1, Z2)]
$\mu_j < 0$ je \emph{rezultat}, ne greška --- dokaz da ta kombinacija aktivnosti
nije optimalna. Tako je odbačen i vrh dopustivog skupa u Z1 ($\mu_2 = -2{,}45$).
\end{zamka}

\begin{zamka}[4. Simulacija ili uzorak predstavljeni kao dokaz (Z2, Z4)]
$0/100$ realizacija bez kršenja nije dokaz robusnosti --- dokaz su vršne
nejednakosti. Rastuća simulacija nije dokaz nestabilnosti --- dokaz je spektralni
radijus monodromijske matrice. Uvijek razdvoji \emph{ilustraciju} od
\emph{dokaza}.
\end{zamka}

\begin{zamka}[5. Predznaci pri prebacivanju (svi zadaci)]
$-x_1-x_2 \leq -10 \Leftrightarrow -x_1-x_2+10 \leq 0$ (Z1);
$a_{11} \in [-1{,}15,\,-0{,}85]$, ne obrnuto (Z2); umnožak izvandijagonalnih
elemenata je $-\pi^2/4$ i u karakterističnu jednadžbu ulazi s minusom (Z4).
Kontrola je uvijek ista: uvrsti jednu poznatu točku i provjeri predznak.
\end{zamka}

\section{Što reći kad ne znaš}
\label{sec:pod-ne-znam}

\begin{izgovor}[Tri rečenice koje uvijek prolaze]
\begin{itemize}[leftmargin=*]
  \item \textbf{Kad ne znaš odgovor:} ,,To nisam računao u radu. Mogu reći kako
        bih pristupio: \dots`` --- pa navedi metodu. Iskren okvir vrijedi više od
        pogađanja.
  \item \textbf{Kad si nesiguran u pretpostavku:} ,,To zadatak ne propisuje, pa
        sam uzeo \dots{} i to je pretpostavka koju sam naveo u radu.`` Razlika
        između ,,zadano je`` i ,,pretpostavio sam`` je ono što se ocjenjuje.
  \item \textbf{Kad ti se čini da je pitanje zamka:} ponovi definiciju prije
        odgovora. ,,Kvadratna stabilnost znači da postoji \emph{jedna zajednička}
        matrica $P$ za sve modove --- uz to, odgovor je \dots``
\end{itemize}
\end{izgovor}

\section{Zadnja provjera prije izlaska}
\label{sec:pod-checklist}

\begin{predznanje}
\begin{enumerate}[leftmargin=*]
  \item Znaš li napamet $\xopt$, $\muopt$ i $\popt$ Zadatka 1? (Sve ima
        nazivnik $33$.)
  \item Znaš li objasniti zašto vrh dopustivog skupa \emph{nije} rješenje?
  \item Možeš li nacrtati sliku Zadatka 1 iz glave, s dvjema antiparalelnim
        strelicama?
  \item Znaš li u jednom retku dokazati da je supremum po pravokutniku u vrhu?
  \item Znaš li računicu $y(9) \geq 3{,}5$ nasuprot $h(9) = 2{,}70$, koja
        objašnjava cijelu sliku Zadatka 3?
  \item Znaš li zašto uvjet $(A+BK)\trans P + P(A+BK) \nd 0$ nije LMI, i što se s
        tim radi?
  \item Znaš li reći $\rho(\Phi_{\mathrm{cikl}}) = 9e^{-2} = 1{,}218 > 1$ i
        objasniti odakle devetka?
  \item Imaš li spremnu rečenicu za svaku pretpostavku koju si sam donio?
\end{enumerate}
\end{predznanje}

\begin{izgovor}[Zadnja napomena]
Ispitivač ne traži da si sve zapamtio. Traži da znaš \emph{odakle} što dolazi ---
je li to bilo u zadatku, na predavanju, ili tvoja odluka. Cijeli je ovaj udžbenik
napisan oko te razlike: siva boja za zadatak, crvena za predavanje, plava za tvoj
rad. Ako na svako pitanje možeš reći kojom je bojom obojano, obrana je gotova.
\end{izgovor}
